\documentclass[11pt]{article}
\usepackage[margin=1in]{geometry}
\usepackage{amsmath,amssymb,amsthm,thm-restate}
\usepackage[hidelinks]{hyperref}
\usepackage[capitalize]{cleveref}
\usepackage{mathtools}
\usepackage{mathrsfs}
\usepackage{microtype}
\usepackage[T1]{fontenc}
\usepackage[utf8]{inputenc} 



\newtheorem{theorem}{Theorem}
\newtheorem{corollary}[theorem]{Corollary}
\newtheorem{conjecture}{Conjecture}
\newtheorem{lemma}[theorem]{Lemma}
\newtheorem{claim}[theorem]{Claim}
\newtheorem{observation}[theorem]{Observation}
\theoremstyle{definition}
\newtheorem{definition}{Definition}


\usepackage{xcolor}
\usepackage[textsize=tiny]{todonotes}
\newcommand{\CC}{{\mathscr C}}
\newcommand{\DD}{{\mathscr D}}
\newcommand{\BB}{{\mathscr B}}
\newcommand{\QQ}{{\mathcal Q}}
\newcommand{\PP}{{\mathcal P}}
\newcommand{\N}{\mathbb{N}}
\newcommand{\jan}[1]{\begin{janin}#1\end{janin}}
% \newcommand{\janin}[1]{\todo[inline,size=\normalsize,color=orange!40,caption={}]{Jan: #1}}
\newcommand{\niko}[1]{\todo[color=green!40]{Niko: #1}}
\newcommand{\nikoin}[1]{\todo[inline,size=\normalsize,color=green!40,caption={}]{Niko: #1}}
\newcommand{\szym}[1]{\todo[color=blue!40]{Szymon: #1}}
\newcommand{\szymin}[1]{\begin{szin}#1\end{szin}}
\newcommand{\mipi}[1]{\todo[color=pink!40]{Michał: #1}}
\newcommand{\mipiin}[1]{\todo[inline,size=\normalsize,color=pink!40,caption={}]{Michał: #1}}
\newcommand{\rose}[1]{\todo[color=purple!40]{Rose: #1}}
\newcommand{\RMM}[1]{\begin{RMMin}#1\end{RMMin}}

\newcommand{\eps}{\varepsilon}
\newcommand{\Oh}{\mathcal{O}}
\newcommand{\inAppendix}{$\clubsuit$}%RMM ADDED
\newcommand{\inducedSS}[2]{#1/#2}%RMM ADDED

\newcommand{\nc}{\text{nc}}
\newcommand{\tc}{\text{tc}}
\renewcommand{\leq}{\leqslant}
\renewcommand{\le}{\leqslant}
\renewcommand{\geq}{\geqslant}
\renewcommand{\ge}{\geqslant}
\renewcommand{\subset}{\subseteq}
\newcommand{\set}[1]{\{#1\}}
\newcommand{\setof}[2]{\set{#1 \,|\, #2 }}
\newcommand{\from}{\colon}
\newcommand{\res}[1]{{\downharpoonright}_{#1}}
\renewcommand{\cal}{\mathcal}
\DeclareMathOperator{\VCdim}{VCdim}
\DeclareMathOperator{\polylog}{polylog}
\DeclareMathOperator{\poly}{poly}
\DeclareMathOperator{\subpoly}{subpoly}

% More compact itemizes
\usepackage{enumitem}
\setlist[itemize]{topsep=4pt,itemsep=3pt,parsep=0pt} 
\setlist[enumerate]{topsep=4pt,itemsep=3pt,parsep=0pt}
\setlist[description]{topsep=4pt,itemsep=3pt,parsep=0pt} 

\usepackage{comment}
\usepackage{tcolorbox}
\tcbuselibrary{skins,breakable}

\newif\ifcomments
\commentstrue        % set to \commentsfalse to hide all comments

\newlength{\NormalParindent}
\newlength{\NormalParskip}
\AtBeginDocument{%
  \setlength{\NormalParindent}{\parindent}%
  \setlength{\NormalParskip}{\parskip}%
}

\newcommand{\createCommentEnv}[3]{%
  \ifcomments
     \newtcolorbox{#1}{
      breakable, enhanced,
      colback=#3!10,
      colframe=#3,
      boxrule=1pt,
      borderline west={2pt}{0pt}{#3},
      left=10pt,right=10pt,top=5pt,bottom=5pt,
      boxsep=0pt,      
      before skip=\NormalParskip,
      after skip=\NormalParskip,      
      title={\bfseries #2:},
      % stick out into margins:
      width=\dimexpr\linewidth+10pt\relax,
      left skip=-10pt,
      % keep normal paragraph formatting:
      before upper={
        \setlength{\parindent}{\NormalParindent}
        \setlength{\parskip}{\NormalParskip}
      },
    }%
  \else
    \excludecomment{#1}%
  \fi
}


\createCommentEnv{szin}{Sz}{red}
\createCommentEnv{janin}{Jan}{green}
\createCommentEnv{RMMin}{RMM}{orange}
\createCommentEnv{claudein}{Claude}{cyan}
\newcommand{\claude}[1]{\begin{claudein}#1\end{claudein}}
\createCommentEnv{oldin}{Old}{gray}
\newcommand{\old}[1]{\begin{oldin}#1\end{oldin}}
\createCommentEnv{newin}{New}{green!50!black}
\newcommand{\new}[1]{\begin{newin}#1\end{newin}}


\let\phi\varphi
\let\epsilon\varepsilon
\newcommand{\Pp}{\mathcal{P}}
\newcommand{\Qq}{\mathcal{Q}}
\newcommand{\Ll}{\mathcal L}
\newcommand{\Dom}{\text{Dom}}
\newcommand{\remph}[1]{{\color{red} \emph{#1}}}
\newcommand{\Aug}{\text{Aug}}
\newcommand{\tf}{\leqslant_{\text{tf}}}
\newcommand{\Tcl}{\text{T-cl}}
\newcommand{\hamming}{\mathrm{Hamming}}
\newcommand{\Graphs}{\mathrm{Graphs}}
\newcommand{\twin}{\mathtt{twin}}

% %%% LEFTBAR BEGIN %%%%%%%%%%%%%%%%%%%%%%%%%%%%%%%%%%%%%%%%%%%%%%%%%%%%%%%%%%%%%
% \usepackage{framed}
% %https://tex.stackexchange.com/questions/50612/alignment-in-the-leftbar-environment-is-rightbar-possible
% \newlength{\leftbarwidth}
% \setlength{\leftbarwidth}{3pt}
% \newlength{\leftbarsep}
% \setlength{\leftbarsep}{10pt}

% \renewenvironment{leftbar}[1][blue]
% {%
% \def\FrameCommand
% {%
% {\hspace{-8pt} \color{#1} \vrule width 3pt}%
% \hspace{0pt}%must no space.
% \fboxsep=\FrameSep\colorbox{#1!5!}%
% }%
% \MakeFramed{\hsize\hsize\advance\hsize-\width\FrameRestore}%
% }
% {\endMakeFramed}
% \setlength{\FrameSep}{5pt}
% %%% LEFTBAR END %%%%%%%%%%%%%%%%%%%%%%%%%%%%%%%%%%%%%%%%%%%%%%%%%%%%%%%%%%%%%%%


% \makeatletter
% %define \szin environment, which which ignores its content if todonotes is used with "disabled" option, and shows a comment otherwise

% \newcommand{\createCommentEnv}[2]{%
%   \if@todonotes@disabled
%     \excludecomment{#1}%
%   \else
%     \newenvironment{#1}{\begin{leftbar}[#2]{\noindent\bf}}{\end{leftbar}}%
%   \fi
% }
% \makeatother

% \createCommentEnv{szin}{red}
% \createCommentEnv{janin}{blue}
% \createCommentEnv{RMM}{orange}

%   \if@todonotes@disabled % Check todnotes is used with disabled option
%   \excludecomment{szin}
%   \excludecomment{janin}
%   \excludecomment{RMM}
%   \else %todonotes enabled
%   \newenvironment{szin}{\begin{leftbar}[red]{\noindent\bf sz: }}{\end{leftbar}}
%   \newenvironment{RMM}{\begin{leftbar}[orang]{\noindent\bf sz: }}{\end{leftbar}}
%   \fi
%   \if@todonotes@disabled % Check todnotes is used with disabled option
%   \excludecomment{janin}
%   \else %todonotes enabled
%   \newenvironment{janin}{\begin{leftbar}[blue]{\noindent\bf Jan: }}{\end{leftbar}}
%   \fi

% \makeatother

\begin{document}

\title{Neighborhood Complexity and Dense Analogue of Bounded Expansion}
\date{}
\maketitle

\section{Definitions}
All classes in this paper are assumed to be \remph{hereditary}, that is, close under taking induced substructures.
A class of graphs $\CC$ is \remph{weakly sparse} there is an integer $t$ such that no graph of $\CC$ contains $K_{t, t}$ as a subgraph. It is \remph{degenerate} if there exists a constant $c$ such that for any graph $G \in \CC$, it holds that $|E(G)|/|V(G)| \le c$.
The language of graphs is denoted by \remph{$\Graphs$}.

\subsection{Transductions}
Given two languages $\Ll$ and $\Ll'$,
an \remph{interpretation from $\Ll$ to $\Ll'$} is
a family $(\varphi_r)_{r \in \text{relations}(\Ll')}$ of $\Ll$-formulae where $\varphi_r$ has as many free variables as the arity of $r$.
If $S$ is a $\Ll$-structure, then \remph{$I(S)$} is the $\Ll'$-structure on $\Dom(S)$ by taking, for each relational symbol $r$ of $\Ll'$, the set of tuples
$\{\overline x \in \Dom(S)~:~S \models \varphi_r(\overline x)\}$. 

Given an interpretation $I$ from $\Ll$ to $\Ll'$,
and an interpretation $J$ from $\Ll'$ to $\Ll''$, the \remph{composition $J \circ I$} is the interpretation from $\Ll$ to $\Ll''$ obtained by substituting, in each formula of $J$, each relation symbol $r$ with the associated formula $\varphi_r$ of $I$.

The \remph{$k$-monadic lift} of a language $\Ll$ is the language
$\Ll^k := \Ll \cup \{U_1, \dots, U_k\}$ where each $U_i$ is a new symbol of unary relation.
An $\Ll$-structure $S$ naturally gives rise to a set \remph{$\Aug_k(S)$} of $\Ll^k$-structures by adding all the possible $k$-unary relations on $S$.

The \remph{$c$-copy} of a language $L$, denoted $(\Ll + c)$ is the language obtained by adding a single $c$-ary relation $R_c$ to $\Ll$.
An $\Ll$-structure $S$ translates into the $(\Ll + c)$-structure $(S + c)$ obtained by taking the disjoint union of $c$ copies of $S$, and letting $R_c$ consist of all tuples $(x_1, \ldots, x_c)$ formed by the $c$ copies of a single element of $S$.

Finally, for two languages $\Ll$ and $\Ll'$, two integers $c$ and $k$, and an interpretation $I$ from $(\Ll + c)^k$ to $\Ll'$, a \remph{$(c, k, I)$-transduction} $T$ is the operation that takes as input an $\Ll$ structure $S$ and returns $\text{Her}(I(\Aug_k(S + c)))$ the hereditary closure of $I(\Aug_k(S + c))$.

A class of $\Ll'$ structures $\DD$ is \remph{transducible from} a class of $\Ll$ structures $\CC$ when there exists a $(c, k, I)$-transduction $T$ such that $\DD \subseteq  T(\CC) := \cup_{S \in \CC} T(S)$.
We will denote this by $\DD \tf \CC$, and it can be observed that $\tf$ is indeed a quasi order.

The \remph{complexity} of a $(c, k, I)$-transduction $T$ is the maximum among $c$, $k$, and the quantifier rank of $I$.
When a transduction has complexity $\le k$, we call it a \remph{$k$-transduction}.
It is well-known that the number of different $k$-transductions on a finite language is finite up to equivalence (depending on $k$).


Given a class $\CC$, we define the \remph{closure by transduction} of $\CC$, denoted by $\Tcl(\CC)$, as the set of classes $\DD$ with $\DD \tf \CC$.

Given a transduction $T$, a graph $G$, and $H \in T(G)$, each vertex of $H$ is a copy of some vertex of $G$.
The set of vertices $v$ of $G$ such that at least one copy of $v$ belongs to $H$ is called the \remph{$T$-support} of $H$ in $G$.
An induced subgraph $G'$ of $G$ is a \remph{$T$-model} of $H$ in $G$ if $V(G')$ contains the $T$-support of $H$ in $G$, and $H \in T(G')$ as a labelled graph.

A transduction $T$ of a class $\CC$ is \remph{linear} when there exists a constant $c := c(\CC, T) > 0$ such that, for any $G \in \CC$, for any $H \in T(G)$, there is a model $G'$ of $H$ in $G$ such that $|V(G')| \le c \cdot |V(H)|$.

\begin{lemma}\label{lem:transunion}
For any two transductions $T_1$ and $T_2$, each of complexity at most $q$, there is a transduction $T$ of complexity $\Oh_q(1)$ such that for any class $\CC$ we have $T_1(\CC) \cup T_2(\CC) \subseteq T(\CC)$.
Moreover, if $T_1$ and $T_2$ are linear on $\CC$, then so is $T$.
\end{lemma}
\begin{proof}
Let $(c_i, k_i, I_i)$ describe the transduction $T_i$.
Then, by augmenting the input structure by $c_1 + c_2$ unary predicates, one can write formulas that apply $I_1$ whenever only the $c_1$ first unary predicates are non-empty, and $I_2$ whenever only the $c_2$ last unary predicates are non-empty.
\end{proof}

\begin{definition}
A class $\CC$ is \remph{monadically dependent} if $\mathcal G \not \tf \CC$ where $\mathcal G$ denotes the class of all graphs.
\end{definition}

\begin{definition}
The Dense Analogue of Bounded Expansion (DABE) is the set of classes $\CC$ such that, for any non degenerate weakly sparse class $\DD$, we have $\DD \not \tf \CC$.
The size-preserving DABE (SP-DABE) is the set of classes that do not transduce a non degenerate weakly sparse class $\DD$ by size-preserving transduction.
\end{definition}

It is clear that a class in DABE is monadically dependent. But in addition, thanks to [TODO ref] one can show that a class in SP-DABE is also monadically dependent.

It is known that the classes that are both weakly sparse and in DABE are the classes of bounded expansion. [TODO ref]

\subsection{Type complexity}
The \remph{neighborhood complexity} of a class of graphs $\CC$ is the function $\nc_{\CC} : \N \to \N$ defined by:
\[
\nc_{\CC}(n) := \max_{G \in \CC} \max_{A \subseteq V(G), |A| \le n} |\{ N(v) \cap A~|~v \in V(G)\}|
\]
For any integer $k$, we define the \remph{k-type complexity }$\tc_{\CC}(k, \cdot) : \N \to \N$ as the maximum among the neighborhood complexities of the graph classes obtained from $\CC$ by application of a $k$-transduction.

\[
\tc_{\CC}(k, n) := \max_{T \text{ a $k$-transduction}} \nc_{T(\CC)}(n)
\]

Note that the finiteness of the set of $k$-transductions allows $\tc$ to always be well-defined.
We say that the class $\CC$ has \remph{linearity} $c$ if $\nc_{\CC}(n) \le c \cdot n$ for any positive $n$. Similarly, it has \remph{$k$-translinearity} $c$ if $\tc_{\CC}(k, n) \le c \cdot n$ for any positive $n$.

A class of graphs is \remph{translinear} when its $k$-translinearity is finite for any $k$.
The main goal of this paper is to prove the following theorem.

\begin{theorem}\label{thm:main}
  A class $\CC$ is translinear if and only if it is in DABE.
\end{theorem}

In addition, we prove the more precise following statement.

\begin{theorem}\label{thm:main}
  Any class in SP-DABE has linear neighborhood complexity.
\end{theorem}

Note that the implication from translinearity to DABE in \Cref{thm:main} follows from a combination of the following known theorems.

\begin{theorem}\label{thm:bddexp2nc}[TODO ref]
  If a weakly sparse class of graphs has bounded expansion, then it has linear neighborhood complexity.
\end{theorem}

\begin{theorem}\label{thm:sparsenc2bddexp}[TODO ref]
  Any translinear, weakly sparse, class of graphs has bounded expansion.
\end{theorem}

\begin{lemma}\label{lem:dreier-sampling-lin}
Let $G = (A, B, E)$ be a bipartite graph with no isolated vertices in $B$, and assume that the average degree of vertices of $B$ is at most $d$.
Then there are $A' \subseteq A$ and $B' \subseteq B$ with $|B'| \ge c_{\ref{lem:dreier-sampling-lin}}(d) \cdot |B|$ such that each vertex in $B'$ has exactly one neighbor in $A'$.
\end{lemma}
\begin{proof}
Since the average degree of vertices in $B$ is at most $d$, at least $|B|/2$ vertices in $B$ have degree at most $2d$.
Let $B_1 \subset B$ be the set of vertices in $B$ of degree at most $2d$.

Let $X$ be a subset of $A$ taken at random, where each vertex $a \in A$ belongs to $X$ with probability $p :=(2d)^{-1}$.
For each $b \in B_1$, let $E_b$ be the event ``$b$ has exactly one neighbor in $X$''.
The vertices of $B$ being non-isolated, the probability $\mathbb{P}(E_b)$ is $\text{deg}(b) \cdot p \cdot (1 - p)^{\text{deg}(b) - 1}$ with $\text{deg}(b)$ the degree of $b$, hence $\mathbb{P}(E_b) \ge c_d := p^{2d}$.
By linearity of expectation, the expected number of vertices of $B_1$ having exactly one neighbor in $X$ is at least $c_d \cdot |B_1|$.
Hence, some outcome $A' \subset A$ achieves at least $c_d \cdot |B_1| \ge c_d/2 \cdot |B|$ such vertices.
It suffices then to take $B' := \{b \in B_1 ~:~ |N(b) \cap A'| = 1\}$.
\end{proof}

\section{Proof of the key lemma}

Given an integer $k$, we call \remph{$k$-equivalence} a formula $\phi(x, y)$ of quantifier rank at most $k$ on the language $\Graphs^k$ that describe an equivalence relation.
Given a $k$-equivalence $\phi(x, y)$ and a graph $G$, we say that a partition $\mathcal P$ of $V(G)$ is \remph{$k, \phi$-definable} if there exists $G^+ \in \Aug_k(G)$ such that $G^+ \models \phi(u, v)$ if and only if $u$ and $v$ are in the same part of $\mathcal P$.
The partition is said to be \remph{$k$-definable} if it is $k, \phi$-definable for some formula $\phi$.
If $\Pp$ is a partition of a subset $S \subseteq V(G)$, we say that it is \remph{$k$-definable in $G$} if there is a $k$-definable partition of $V(G)$ that equals $\Pp$ when restricted to $S$.

Given a graph $G$ and two partitions $\mathcal P$ and $\mathcal Q$ that are $k$-definable in $G$, their \remph{intersection graph} $\text{Inter}(\mathcal P, \mathcal Q)$ is the graph with vertex set $\{p \in \mathcal P\} \uplus \{q \in \mathcal Q\} \uplus \{(p,q) ~:~ p \in \mathcal P, q \in \mathcal Q, p \cap q \neq \emptyset\}$ and edges joining each vertex $(p, q)$ to the two vertices $p \in \mathcal P$ and $q \in \mathcal Q$.
The set of vertices $\{(p,q) ~:~ p \in \mathcal P, q \in \mathcal Q, p \cap q \neq \emptyset\}$ is the \remph{common refinement of $\mathcal P$ and $\mathcal Q$} denoted $\mathcal P \wedge \mathcal Q$.

A $k$-equivalence $\phi$ is \remph{size-preserving} in the class $\CC$ whenever, for any graph $G \in \CC$, for any $A \subseteq V(G)$, any $k, \phi$-definable partition of $A$ is already $k, \phi$-definable in $G[\overline{A}]$ for some $A \subseteq \overline{A}$ with $|\overline{A}| \le \Oh_{\phi}(|A|)$.

\begin{lemma}\label{lem:twin_SP}
The twin relation over a set $S$, $\twin(x, y) := \forall s, S(s) \Rightarrow \left(E(x, s) \Leftrightarrow E(y, s) \right)$ is size preserving.
\end{lemma}
\begin{proof}
Let $G$ be a graph, and $A \subseteq V(G)$.
Let $S \subseteq V(G)$ and $\Pp$ be the partition of $A$ into twin classes over $S$.
We prove that there exists a set $S_A \subseteq S$ of size at most $|A - 1|$ such that two vertices of $A$ are twins over $S$ if and only if they are twins over $S_A$. Since this implies that $\Pp$ is $\twin$-definable in $G[A \cup S_A]$, it is sufficient to prove that $\twin$ is size-preserving.

We construct the set $S_A \subseteq S$ inductively.
We start setting $S_A := \emptyset$, and we keep track of the size of the partition $\Pp_{S_A}$ of $A$ into twin classes over $S_A$.
We will ensure that $S_A \subseteq S$, so the partition $\Pp_{S_A}$ will be a coarsening of $\Pp$, hence we terminate if $|\Pp_{S_A}| = |\Pp|$.
Otherwise, there are two elements $a, a' \in A$ that are twins over $S_A$ but not over $S$. In particular, there exists $s \in S$ such that $E(s, a) \Leftrightarrow \neg E(s, a')$.
By adding $s$ to $S_A$, we increase the size of $S_A$ by one, and the size of $\Pp_{S_A}$ by at least one.

\end{proof}
\begin{lemma}\label{lem:transInter}
    Let $\phi, \psi$ be two $k$-equivalences, the there exists a transduction $T_{\phi, \psi}$ of complexity $\Oh_k(1)$ such that,
    for any class $\CC$ and any graph $G \in \CC$, for any $S \subseteq V(G)$ and any two partitions $\mathcal P$ and $\mathcal Q$ of $S$, respectively $\phi$ and $\psi$-definable, the graph $\text{Inter}(\Pp , \Qq)$ is in $T_{\phi, \psi}(G)$.
    In addition, if both $\phi$ and $\psi$ are size preserving, then so is $T_{\phi, \psi}$.
\end{lemma}
\begin{proof}

Let $U_1^\phi, \ldots, U_k^\phi$ be the $k$ unary predicate $\phi$ may use; and let $U_j^\psi, j \in [k]$ be the analogous data for $\psi$. 

We consider the $3$-copy structure $(G+3)$ with copies $V^{(1)}, V^{(2)}, V^{(3)}$, with the relation $\text{copy}(x, y)$ asserting that $x$ and $y$ are copies of a same vertex of $V(G)$. Then, we augment it as follows: add the unary predicates $U_j^\phi$ and $U_j^\psi$ for each $j \in [k]$ on the first copy; a predicate $\mathrm{Rep}_{\Pp}$ picking one representative per $\Pp$-class intersecting $S$ in $V^{(1)}$; a predicate $\mathrm{Rep}_{\Qq}$ picking one representative per $\Qq$-class intersecting $S$ in $V^{(2)}$; and a predicate $\mathrm{Rep}_{\Pp \Qq}$ picking a representative in $V^{(3)}$ for each non-empty $p \cap q \cap S$ with $p \in \Pp$ and $q \in \Qq$.

It remains to define the edge relation.
In order to do this, it suffices to add edges $xy$ that satisfy expressions of the form
\[
\theta_E(x, y) := \mathrm{Rep}_{\Pp}(x) \wedge \mathrm{Rep}_{\Pp \Qq}(y) \wedge \exists y', V^{(1)}(y') \wedge \text{copy}(y, y') \wedge  \phi_{\Pp}(x, y')
\]
and the symmetric variants.
The graph $Inter(\Pp, \Qq)$ is can be found in $T_{\phi, \psi}(G)$, as the graph induced by $\{v \in  \mathrm{Rep}_{\Pp} \cup \mathrm{Rep}_{\Qq} \cup \mathrm{Rep}_{\Pp \wedge \Qq}\}$.

Now, assume that $\phi$ and $\psi$ are size-preserving. Then, once $\Pp$ and $\Qq$ are fixed, one can find a subset $\overline{S}$ containing $S$ such that $\Pp$ and $\Qq$ are respectively $\phi$ and $\psi$ definable in $G[\overline{S}]$, and $|\overline{S}| \le \Oh_{\phi, \psi}(|S|)$.
By the previous proof, we obtain that $Inter(\Pp, \Qq) \in T_{\phi, \psi}(G[\overline{S}])$.
\end{proof}

This will allow us to prove the following central tool of this paper.

\begin{corollary}\label{cor:2equiv_ref}
    Let $\CC$ be a class in DABE (resp. in SP-DABE).
    Let $G \in \CC$, $S \subseteq V(G)$ and $\Pp$ and $\Qq$ be any two partitions of $S$ that are $k$-definable (resp. by size preserving $k$-equivalences) in $G$.
    Then there is a constant $c_{\ref{cor:2equiv_ref}}$ depending only on $\CC$ and $k$ such that $|\Qq \wedge \Pp| \le c_{\ref{cor:2equiv_ref}} \cdot (|\Pp| + |\Qq|)$.
\end{corollary}
\begin{proof}
Let $\mathcal F_k$ be a maximal set of non-equivalent (as logical statement) $k$-equivalences. The set $\mathcal F_k$ is finite.
For any $\phi, \psi \in \mathcal F_k$, we define the class
\[
\DD_{\phi, \psi} :=
\{\text{Inter}(\Pp,\Qq) ~:~ G \in \CC \text{, and $\Pp$ and $\Qq$ respectively $\phi$ and $\psi$-definable partitions of } S\}.
\]
By \Cref{lem:transInter}, $\DD_{\phi, \psi}$ is $c$-transducible from $G$ for some $c$ depending only on $k$, and so by \Cref{lem:transunion} $\DD_k := \bigcup_{\phi, \psi \in \mathcal F_k} \DD_{\phi, \psi}$ is also transducible from $\CC$ by a transduction whose complexity depends only on $k$.
Note that for any set $S \subseteq V(G)$ if $\Pp_S$ and $\Qq_S$ are two partitions of $S$ that are $k$-definable in $G$, then $\text{Inter}(\Pp_S, \Qq_S)$ is in $\DD_k$.

For any graph $G$ and any $k$-definable partitions $\Pp, \Qq$ of $V(G)$, the graph $\text{Inter}(\Pp, \Qq)$ is bipartite between $A := \Pp \wedge \Qq$ and $B := \Pp \uplus \Qq$.
In addition, every vertex of $A$ has degree exactly~$2$.
Hence, $\DD_k$ is a class of weakly sparse graphs that is transducible from $\CC$.
Since the later is in DABE, we deduce that $\DD_k$ is a class of bounded expansion.
By \Cref{thm:bddexp2nc} classes of bounded expansion have linear neighborhood complexity, hence there exists a constant $c$ such that $\DD_k$ has linearity at most $c$.

Finally, for any two $k$-definable partitions $\Pp, \Qq$, the definition of $\text{Inter}(\Pp, \Qq)$ ensures that two different vertices in $\Pp \wedge \Qq$ have distinct neighborhoods over $\Pp \uplus \Qq$. Hence the linearity of $\DD_k$ ensures that $|\Pp \wedge \Qq| \le c \cdot (|\Pp| + |\Qq|)$ for some constant $c$, completing the proof by taking $c_{\ref{cor:2equiv_ref}} := c$.


The proof for $\CC$ is in SP-DABE is exactly the same, apart from the fact that the classes $\DD_{\phi, \psi}$ are defined only for $\phi$ and $\psi$ size preserving $k$-equivalences.
\end{proof}

\section{Linear Neighborhood Complexity for DABE Classes}
\label{sec:nc-dabe}

In this section we prove that every class in DABE has \emph{linear} neighborhood complexity.

\subsection{Preliminaries}

Any bipartite graph $G=(A,B,E)$ induces a \remph{set system} $\inducedSS{B}{A} := \setof{N(b)\cap A}{b\in B}$ on $A$.
For $A'\subseteq A$ and $B'\subseteq B$ we write $\inducedSS{B'}{A'}$ for the set system $\setof{N(b)\cap A'}{b\in B'}$ on $A'$. Given a set-system $\cal F$ on $V$, and $X \subseteq V$ we will write \remph{$\cal F\res X$} for its restriction to $X$.
Two vertices of $B$ are \remph{twins} if they have equal neighborhoods in $A$; $|\inducedSS{B}{A}|=|B|$ iff $B$ contains no twins.

Given a set system $\cal F$ on $V$, a set $X \subseteq V$ is \remph{shattered} if $\cal F\res X=2^X$.
The \remph{VC-dimension} of a set system $\cal F$ on $V$ is the size of a largest shattered set (and $-\infty$ if $\cal F=\emptyset$).
We write $\VCdim(G)$ for $\VCdim(\inducedSS{V(G)}{V(G)})$. The VC-dimension of a class of graphs $\CC$ is $\sup_{G \in \CC} \VCdim(G)$.
Let $v\in V$. Sets $F,F' \in \cal F$ form a \remph{$v$-pair} if $F\triangle F'=\set{v}$.
The \remph{Hamming graph} $\hamming(\cal F)$ of a set system $\cal F$ on $V$ is the graph with vertex set $\cal F$ in which $F$ and $F'$ are adjacent whenever $|F \triangle F'| = 1$, i.e. whenever they form a $v$-pair for some $v \in V$.

If $F$ belongs to some $v$-pair it is \remph{$v$-positive} if $v\in F$, \remph{$v$-negative} otherwise.
Given $G=(A,B,E)$, the \remph{positive} (resp.\ \remph{negative}) \remph{merge graph} of $G$ is the bipartite graph with parts $A$ and $\inducedSS{B}{A}$ where $a\in A$ is adjacent to $F\in \inducedSS{B}{A}$ iff $F$ is $a$-positive (resp.\ $a$-negative).
The \remph{merge graph} of $G$ is the union (intended as edge union) of its positive and negative merge graphs.

\begin{lemma}\label{lem:dep-bounded-vc}
Every monadically dependent class has bounded VC-dimension.
\end{lemma}

\begin{lemma}[TODO ref]\label{lem:VC-dim2deg}
The Hamming graph of a set-system $\cal F$ of VC-dimension at most $d$ has at most $d |\cal F|$ edges.  
\end{lemma}

\begin{corollary}\label{cor:hamming}
Let $\cal F$ be a set-system of VC-dimension $d$, whose Hamming graph $H$ has at least $m$ edges. Then $H$ has at least $m/d$ non-isolated vertices.
\end{corollary}
\begin{proof}
Indeed, let $\cal F'$ be the set of non-isolated vertices of $H$. We have that $|E(H[\cal F'])| = |E(H)| \ge m$, and the bound on the VC-dimension together with \Cref{lem:VC-dim2deg} gives that $|\cal F'| \ge m / d$.
\end{proof}

\begin{observation}\label{obs:merge-graph-density}
Let $G = (A, B, E)$ be a bipartite graph of VC-dimension at most $d$ and without twins in $B$, and let $B' \subseteq B$ the set of non-isolated vertices of the merge graph of $G$.
Then the average degree of vertices of $B'$ in the merge graph of $G$ is at most $d$.
\end{observation}

\begin{lemma}[VC-dimension induction]\label{lem:VC-induction-dabe}
For any nonempty set system $\cal F$ on $V$ and $v\in V$, the two set systems
\[
\cal F^+ := \{F \in \cal F ~|~F \text{ is $v$-positive}\} \text{ and } \cal F^- := \{F \in \cal F ~|~F \text{ is $v$-negative}\}
\]
have VC-dimension strictly less than $\VCdim(\cal F)$.
\end{lemma}
\begin{proof}
Suppose $X \subseteq V$ is shattered by $\cal F^+$. Every member of $\cal F^+$
contains $v$, so $v \notin X$. We claim $X \cup \{v\}$ is shattered by $\cal F$.
Let $Y \subseteq X$. Pick $F \in \cal F^+$ with $F \cap X = Y$. Then
$F \cap (X \cup \set{v}) = Y \cup \set{v}$, and since $F$ is $v$-positive,
$F \setminus \set{v} \in \cal F$ and
$(F \setminus \set v) \cap (X \cup \set v) = Y$. Hence every subset of
$X \cup \set v$ is realised, so
$\VCdim(\cal F) \ge |X| + 1 > \VCdim(\cal F^+)$. The negative case is symmetric.
\end{proof}

Given a class of graphs $\CC$, we denote by \remph{$\cal M^+_{\CC}$}(resp. \remph{$\cal M^-_{\CC}$}, \remph{$\cal M_{\CC}$}) the class of the positive merge graphs of the bipartite graphs of $\CC$ (resp. negative merge graphs, merge graphs).

\subsection{Definable partitions and families of set systems}
When $\Pp$ is a partition of $B$, then $(G, \Pp)$ induces the \remph{family of set systems} $(p/A)_{p \in \Pp}$ on $A$ that contains one set system for each part of $\Pp$. We say that the family is \remph{$k$-definable} or whenever the partition is.

\begin{lemma}\label{lem:partition_size}
Let $\CC$ be in SP-DABE and $k$ be an integer. Then there is a constant $c_{\ref{lem:partition_size}}(\CC, k) $ such that
for any size preserving $k$-equivalence $\phi$, for any bipartite graph $G = (A, B, E) \in \CC$, any $A_0 \subseteq A, B_0 \subseteq B$ and any $\phi$-definable partition $\Pp$ of $B_0$, we have $\sum_{p \in \Pp} |p/A_0| \leq c_{\ref{lem:partition_size}}(\CC, k) \cdot (|\Pp| + |B_0/A_0|)$.
\end{lemma}
\begin{proof}
Let $\Qq$ be the partition of $B_0$ into twin classes over $A_0$, that is, $\Qq$ is defined by $\twin$ over $A_0$. Note that $\Qq$ is in bijection with $B_0/A_0$.

Both partitions $\Pp$ and $\Qq$ are $\max(k, 1) \le k+1$-definable and $\twin$ is size-preserving by \Cref{lem:twin_SP}, so by \Cref{cor:2equiv_ref}, there exists a constant $c_{\ref{lem:partition_size}} := c_{\ref{cor:2equiv_ref}} (\CC, k+1)$ such that 
\[
|\Pp \wedge \Qq| \leqslant c_{\ref{lem:partition_size}} \cdot (|\Pp| + |\Qq|) = c_{\ref{lem:partition_size}} \cdot (|\Pp| + |B/A|).
\]
Finally, for a given $p \in \Pp$, $|p/A_0|$ is the number of classes of $\Qq$ that meet $p$, so $\sum_{p \in \Pp} |p/A_0| = |\Pp \wedge \Qq|$, and the lemma follows.
\end{proof}


\subsection{Distillation procedure}
Our argument will be centered on an induction over the VC-dimension of a $k$-definable family of set systems $(p/A)_{p \in \Pp}$, and will need the total number of Hamming edges $\sum_{p \in \Pp} |E(\hamming(p/A))|$ to be a constant fraction of $\sum_{p \in \Pp} |p/A|$. The next lemma provides this, at the cost of restricting $A$ to a constant fraction of itself.

We say that a bipartite graph $G = (A, B, E) \in \CC$ with $A \neq \emptyset$ has \remph{maximal relative complexity in $\CC$}, or simply \remph{maximal complexity} when the context is clear, if it is twin-free and for any bipartite graph $G' = (A', B', E') \in \CC$ with $0 < |A'| \leq |A|$ we have
\[
\frac{|B'/A'|}{|A'|} \leq \frac{|B/A|}{|A|}.
\]

\begin{lemma}[Distillation procedure]\label{lem:distil}
For any class $\CC$ in SP-DABE and integer $k$, there is a constant
$c_{\ref{lem:distil}}(\CC, k)$ such that for any $\alpha \in (0, 1]$, the following holds.

For any bipartite $G = (A, B, E) \in \CC$ of maximal complexity,
for any $A_0 \subseteq A, B_0 \subseteq B$ such that $|\inducedSS{B_0}{A_0}| \geqslant \alpha \cdot |\inducedSS{B}{A}|$
and any size-preserving $k$-definable partition $\Pp$ of $B_0$ with $|\Pp| \leqslant \frac{|\inducedSS{B_0}{A_0}|}{4 c_{\ref{lem:partition_size}(\CC, k)}}$,
there is an $A_1 \subseteq A_0$, with $|A_1| \geq \alpha \cdot c_{\ref{lem:distil}}(\CC, k) \cdot |A_0|$ such that  
\[
\sum_{P \in \Pp} |E(\hamming(\inducedSS{P}{A_1}))| \ge
\alpha \cdot c_{\ref{lem:distil}}(\CC, k) \cdot \sum_{P \in \Pp} |\inducedSS{P}{A_0}| \ge \alpha c_{\ref{lem:distil}}(\CC, k) \cdot |B_0/A_0|.
\]
\end{lemma}

\begin{proof}
We shrink $A_0$ to $A_1$ iteratively. Set initially $X := A_0$. As long as there is a vertex $v \in X$ such that 
\[
\sum_{P \in \Pp}|\inducedSS{P}{X}| - |\inducedSS{P}{(X - v)}| \leq \frac{|\inducedSS{B_0}{A_0}|}{2|A_0|}
\]
remove $v$ from $X$ and repeat. Otherwise, if there is no such $v$, terminate.

\begin{claim}
At the end of the process $|X| > c \cdot |A_0| $ for 
$c := \frac{\alpha}{8c_{\ref{lem:partition_size}}}$.
\end{claim}
\begin{proof}
By contradiction, assume that eventually the process reaches a non-empty set $X$ of size $\le c \cdot |A_0|$.
Let $\ell := |A_0| - |X|$ and $v_1, \dots, v_{\ell}$ be the vertices of $A_0$ removed iteratively by the process.
We have that
\begin{align*}
\sum_{P \in \Pp}|\inducedSS{P}{X}| =
\sum_{P \in \Pp}|\inducedSS{P}{A_0}|
&-\sum_{P \in \Pp}\left( |\inducedSS{P}{A_0}| - |\inducedSS{P}{A_0 - v_1}|\right) - \dots \\
  &- \sum_{P \in \Pp}\left( |\inducedSS{P}{A_0 - \{v_1, \dots, v_{\ell-1}\}}| - |\inducedSS{P}{A_0 - \{v_1, \dots, v_{\ell}\}}|\right).
\end{align*}
In particular, 
\[
\sum_{P \in \Pp}|\inducedSS{P}{X}| \ge
|B_0/A_0| - \ell \cdot \frac{|\inducedSS{B_0}{A_0}|}{2|A_0|} \geq
\frac{|\inducedSS{B_0}{A_0}|}{2}.
\]
By \Cref{lem:partition_size} applied to $X$ and $B_0$, we know that $\sum_{P \in \Pp}|\inducedSS{P}{X}| \leqslant c_{\ref{lem:partition_size}} \cdot (|\Pp| + |\inducedSS{B_0}{X}|)$.
This gives 
\[
\frac{|\inducedSS{B_0}{A_0}|}{4} \le
\frac{|\inducedSS{B_0}{A_0}|}{2} - c_{\ref{lem:partition_size}} |\Pp| \le
\sum_{P \in \Pp}|\inducedSS{P}{X}| - c_{\ref{lem:partition_size}} |\Pp|
\leqslant c_{\ref{lem:partition_size}} |\inducedSS{B_0}{X}|
\]
We now evaluate the relative complexity of $G[X \cup B_0]$. Since $|X| \le c \cdot |A_0|$, we have
\[
\frac{|\inducedSS{B_0}{X}|}{|X|} \ge
\frac{|\inducedSS{B_0}{A_0}|}{4c_{\ref{lem:partition_size}}} \cdot \frac{1}{|X|} \ge
\frac{\alpha \cdot |\inducedSS{B}{A}|}{4 c_{\ref{lem:partition_size}}} \cdot \frac{1}{\frac{\alpha}{8 c_{\ref{lem:partition_size}}} \cdot |A|} \geqslant \frac{2|\inducedSS{B}{A}|}{|A|}.
\]
Since $G[X \cup B_0] \in \CC$, this contradicts the maximal complexity of $G$.
\end{proof}
Let $A_1$ be the set $X$ at which the process stops.
We have
\[
\sum_{P \in \Pp} |E(\hamming(\inducedSS{P}{A_1}))| =
\sum_{a \in A_1} \sum_{P \in \Pp} |\inducedSS{P}{A_1}| - |\inducedSS{P}{A_1 - a}| \geqslant
\frac{|A_1| \cdot |\inducedSS{B_0}{A_0}|}{2|A_0|} \geqslant
\frac 12 \cdot \frac{\alpha}{8 c_{\ref{lem:partition_size}}}|\inducedSS{B_0}{A_0}|
\]
It remains to fix $c_{\ref{lem:distil}} := \frac{1}{2 \cdot 8c_{\ref{lem:partition_size}}}$ to conclude.
\end{proof}

\subsection{Sparsification}

\begin{definition}
  Let $G=(A,B,E)$ be a bipartite graph and $k\in\N$.
  A \remph{$k$-sparsification} consists of:
  \begin{itemize}
    \item sets $A_0\subseteq A$ and $B_0\subseteq B$, and
    \item functions $f_1,\ldots,f_k\from B_0\to A$,
  \end{itemize}
  such that
  the mapping $b\mapsto (N(b)\cap A_0,f_1(b),\ldots,f_k(b))$ is an injection from $B_0$ to $2^{A_0} \times A^k$.
For such a $k$-sparsification we define:
\begin{itemize}
  \item The \emph{associated partition} $\cal P$ as the partition of $B_0$ such that two vertices $b,b'$ are in the same part of $\cal P$ if and only if 
  $f_j(b)=f_j(b')$ for all $j=1,\ldots,k$. Note that vertices in the same part of $\cal P$ have pairwise distinct neighborhoods in $A_0$;
  \item The \emph{size} as $s=|B_0|$;
  \item The \emph{dimension} as $\max_{P\in\cal P} \VCdim(\inducedSS{P}{A_0})$; and
  \item The \emph{complexity} as the least $q\ge0$ such that there are first-order formulas $\phi_1(x,y),\ldots,\phi_k(x,y)$ of quantifier rank ${\le}q$, each using ${\le}q$ unary predicates, and there is an expansion $G'$ of $G$ with $q$ unary predicates,
  such that $\phi_j$ defines $f_j$ in $G'$ for $j=1,\ldots,k$, so that 
  \[
    f_j(b)=a\iff G'\models \phi_j(b,a)
    \quad
    \text{for all $b\in B_0$ and $a\in A$.}  
  \]
\end{itemize}
When the input graph comes from a DABE class $\CC$, we say that a sparsification is \remph{terminal} if $|B_0/A_0| \le 4 c_{\ref{lem:partition_size}} \cdot|\Pp|$, and
\remph{nonterminal} otherwise.
\end{definition}


Note that any bipartite graph with no twins in $B$ has a trivial $0$-sparsification ($A_0=A$, $B_0=B$, no functions) of size $|B|$, complexity $0$, and dimension $\VCdim(\inducedSS{B}{A})$.

We prove the following, for bipartite graphs $G$ of maximal complexity from a class $\CC$ in DABE:
\begin{enumerate}
  \item a nonterminal $k$-sparsification can be improved to a $(k+1)$-sparsification with strictly smaller dimension, by losing only a constant factor depending only on $\CC$, and increasing the complexity only by a constant (\Cref{lem:step});
  \item repeating this argument, we reach a terminal sparsification after at most $d$ steps, where $d$ upper bounds the VC-dimension of every $G$ in $\CC$ (\Cref{lem:iterate});
  \item  a terminal $k$-sparsification of $G\in \CC$ of complexity $q$ has size $s\le O_{\CC, k, q}(|A|)$ (\Cref{lem:terminal}).
\end{enumerate}
Combining these three points yields a terminal $k$-sparsification with $k\le d$ and size $s$ satisfying
\[
\text{const}(\CC) \cdot |B| \le s \le \text{CONST}(\CC) \cdot |A|
\]
for some constants $\text{const}(\CC) > 0$ and $\text{CONST}(\CC)$.
We now prove the necessary lemmas.

\begin{lemma}\label{lem:terminal}
  Let $\CC$ be a class of bipartite graphs in DABE, and let $k,q\in\N$.
  Let $G=(A,B,E)$ be a graph of $\CC$ admitting a terminal $k$-sparsification of complexity at most $q$ and size $s$.
  Then
  \[
    s\le c_{\ref{lem:terminal}}(\CC, k, q) \cdot |A|.
  \]
\end{lemma}
\begin{proof}
Let $A_0, B_0, f_1, \dots, f_k$ be a terminal $k$-sparsification of complexity at most $q$ and size $s$ of $G$. Let $\Pp$ be the associated partition.
The partition $\Pp$ is $(q+1)$-definable, and we have $|s| = \sum_{P \in \Pp} |\inducedSS{P}{A_0}|$.
Hence, by \cref{lem:partition_size}, we have $c_{\ref{lem:partition_size}}\cdot (|\Pp| + |\inducedSS{B_0}{A_0}|) \ge s$, and by terminality we know $4 c_{\ref{lem:partition_size}} \cdot |\Pp| \ge |\inducedSS{B_0}{A_0}|$. Hence, it suffices to prove that $|\Pp|\le O_{\CC, k, q}(|A|)$.

For each $i \in [k]$ let $\Qq_i$ be the partition of $B_0$ where two vertices $b$ and $b'$ are in the same part if and only if $f_i(b) = f_i(b')$.
By definition, we have $|\Qq_i| \le |A|$.
The partition $\Pp$ is the refinement of the partitions $\Qq_1, \dots, \Qq_k$, that is $\Pp = \Qq_1 \wedge \dots \wedge \Qq_k$.
In addition, since each function $f_i, i \in [k]$ is definable by a formula of quantifier rank at most $q$ using at most $q$ unary predicates, each partition of the form $\bigwedge_{i \in I} \Qq_i$, for $I \subseteq [k]$, is $kq$-definable (one may need $kq$ unary predicates in total).
Hence, by \Cref{cor:2equiv_ref} we obtain the bound $|\Pp| \le O_{k, q, \CC} (|\Qq_1| + \dots + |\Qq_k|) \le O_{k, q, \CC}(|A|)$.
\end{proof}

\section{Main induction}
\begin{lemma}\label{lem:step}
  Let $\CC$ be a class in DABE.
  Suppose $G=(A,B,E) \in \CC$ is of maximal complexity, and has a nonterminal $k$-sparsification of size $s \ge \alpha \cdot |B|$ for $\alpha >0$,
  dimension $d$, and complexity $q$.
  Then there is a $(k+1)$-sparsification of $G$ with:
  \begin{itemize}
    \item size at least $c_{\ref{lem:step}} \cdot |B|$,
    \item dimension at most $d-1$,
    \item complexity at most $q+5$.
  \end{itemize}
  where $c_{\ref{lem:step}} := c_{\ref{lem:step}}(\CC, k, d, q, \alpha)$ is independent of $G$ and non decreasing in $\alpha$.
\end{lemma}
\begin{proof}
Let $A_0\subseteq A$, $B_0\subseteq B$, and $f_1,\ldots,f_k\from B_0\to A$
witness the given $k$-sparsification.
Let $\cal P$ be its associated partition.
For every part $P\in\cal P$, the functions $f_1,\ldots,f_k$ are constant on $P$,
so the injectivity condition in the definition of a sparsification implies that
the map $b\mapsto N_G(b)\cap A_0$ is injective on $P$.
Therefore,
\[
  |\inducedSS{P}{A_0}|=|P|
  \qquad\text{for every }P\in\cal P,
\]
and hence
\[
  \sum_{P\in\cal P}|\inducedSS{P}{A_0}|=\sum_{P\in\cal P}|P|=|B_0|=s.
\]

The sparsification being non-terminal, we have $4 c_{\ref{lem:partition_size}} \cdot |\Pp| \le |\inducedSS{B_0}{A_0}|$, and so we can control the size of $\inducedSS{B_0}{A_0}$ by an application of \cref{lem:partition_size} on the set system $\inducedSS{B_0}{A_0}$ and the partition $\Pp$:
\[s = \sum_{P \in \Pp} |\inducedSS{P}{A_0}| \le c_{\ref{lem:partition_size}} \cdot (|\Pp| + |B_0/A_0|), \text{ hence } |\inducedSS{B_0}{A_0}| \ge \frac{s}{2 c_{\ref{lem:partition_size}}} \ge \frac{\alpha}{2 c_{\ref{lem:partition_size}}} |B| \ge \frac{\alpha}{2 c_{\ref{lem:partition_size}}} |\inducedSS{B}{A}|.
\]
\noindent In addition, we have $d\ge 1$, since otherwise every set system $\inducedSS{P}{A_0}$
would have VC-dimension $0$ and hence size at most $1$, implying $|\Pp| \ge s$ contradicting the non-terminality.
Similarly, we have $|A_0| \ge 1$.

The graph $G$ has maximal complexity, and $\Pp$ is $kq$-definable with $|\Pp| \le \frac{|\inducedSS{B_0}{A_0}|}{4 c_{\ref{lem:partition_size}}}$, so we are able to apply \cref{lem:distil} with $\alpha' := \frac{\alpha}{2 c_{\ref{lem:partition_size}}}$ in place of $\alpha$, on the subsets $A_0 \subseteq A$, $B_0 \subseteq B$ and the partition
$\cal P$ of $B_0$.
This gives a set $A_1\subseteq A_0$ such that
\begin{equation}\label{eq:step-many-hamming-edges}
  \sum_{P\in\cal P}|E(\hamming(\inducedSS{P}{A_1}))|
  \ge h \cdot |\inducedSS{B_0}{A_0}| \text{, with } h := \frac{\alpha}{2 c_{\ref{lem:partition_size}}} c_{\ref{lem:distil}}.
\end{equation}

Choose a set $B^*\subseteq B_0$ so that for each part $P\in\cal P$,
\[
  \inducedSS{(P\cap B^*)}{A_1}=\inducedSS{P}{A_1}
\]
and no two vertices in $P\cap B^*$ have equal neighborhoods in $A_1$.
Thus, the map $b\mapsto N_G(b)\cap A_1$ is a bijection from $P\cap B^*$ to
$\inducedSS{P}{A_1}$.

For each $P\in\cal P$, let $H_P$ be the Hamming graph of
$\inducedSS{P}{A_1}$, with vertex set identified with $P\cap B^*$ (identified using the previously constructed bijection between the two).
Let $H^*$ be the disjoint union of the graphs $H_P$, for $P\in\cal P$.
Then $V(H^*)=B^*$ and, by \eqref{eq:step-many-hamming-edges},
\[
  |E(H^*)|
  =\sum_{P\in\cal P}|E(H_P)|
  \ge h \cdot |\inducedSS{B_0}{A_0}|.
\]

For each $P\in\cal P$, let $M_P^+$ and $M_P^-$ denote the positive and
negative merge graphs of $\inducedSS{P}{A_1}$.
Let $M^+$ and $M^-$ be the disjoint unions of the graphs
$M_P^+$ and $M_P^-$, respectively, identified along the common part $A_1$.
As $A_1\subset A_0$,
% \By \Cref{obs:vc-hereditary},
 we have $\VCdim(\inducedSS{P}{A_1})\le \VCdim(\inducedSS{P}{A_0})\le d$ for every $P\in\cal P$.
Then \Cref{cor:hamming} implies that $H_P$ has at least $|E(H_P)|/d$
non-isolated vertices.
Summing over $P\in\cal P$, we infer that $H^*$ has at least
\[
  \frac{|E(H^*)|}{d}
  \ge \frac{h}{d} \cdot |\inducedSS{B_0}{A_0}|
\]
non-isolated vertices.
Every such vertex is non-isolated in at least one
of $M^+$ and $M^-$, so for some $\sigma\in\set{+,-}$ the graph $M^\sigma$
has at least
\begin{equation}\label{eq:step-many-nonis}
  \frac{h}{2d} \cdot |\inducedSS{B_0}{A_0}|
\end{equation}
non-isolated vertices in $B^*$.

We now choose a large subset on which each vertex has a unique
$M^\sigma$-neighbor.
By \Cref{obs:merge-graph-density}, the non-isolated vertices in $B^*$ graph $M^{\sigma}$ have average degree at most $d$, hence by \cref{lem:dreier-sampling-lin} to the bipartite graph obtained
from $M^\sigma$ by removing the isolated vertices. % from $B^*$.
We obtain sets $A_1'\subseteq A_1$ and $B_1\subseteq B^*$ such that every
vertex of $B_1$ has exactly one neighbor in $A_1'$ in $M^\sigma$, and
\begin{equation}\label{eq:after_sampling}
  |B_1|
  \ge c_{\cref{eq:after_sampling}} \cdot |\inducedSS{B_0}{A_0}|, \text{ with } c_{\cref{eq:after_sampling}} := c_{\ref{lem:dreier-sampling-lin}}
    (d)
    \cdot \frac{h} {2d}.
\end{equation}

Define $f_{k+1}\from B_1\to A_1'$ by mapping each $b\in B_1$ to its unique
neighbor in $A_1'$ in the graph $M^\sigma$.
Let $\cal P_1$ be the family of all nonempty sets of the form
$f_{k+1}^{-1}(v)\cap P$, where $P\in\cal P$ and $v\in A_1'$.

We claim that $A_1$, $B_1$, and the functions
$f_1,\ldots,f_k,f_{k+1}$ form a $(k+1)$-sparsification of $G$.
Indeed, if $b,b'\in B_1$ lie in different parts of $\cal P$, then
$(f_1(b),\ldots,f_k(b))\neq (f_1(b'),\ldots,f_k(b'))$.
If they lie in the same part $P\in\cal P$, then 
$N_G(b)\cap A_1\neq N_G(b')\cap A_1$ whenever $b\neq b'$.
This is because we have chosen $B^* \supseteq B_1$ such that all vertices in $P \cap B^*$ have pairwise different neighborhoods on $A_1$ in $G$.
Hence, the map
\[
  b\mapsto (N_G(b)\cap A_1,f_1(b),\ldots,f_k(b),f_{k+1}(b))
\]
is injective on $B_1$.
Moreover, by construction, $\cal P_1$ is exactly the associated partition of
this new sparsification.

We now verify the claimed properties.
First, let $P'\in\cal P_1$.
Then $P'=f_{k+1}^{-1}(v)\cap P$ for some $P\in\cal P$ and $v\in A_1'$.
The set system $\inducedSS{P'}{A_1}$ is a subfamily of the neighborhood of $v$
in the $\sigma$-merge graph of $\inducedSS{P}{A_1}$, so
\Cref{lem:VC-induction-dabe} yields
\[
  \VCdim(\inducedSS{P'}{A_1})\le \VCdim(\inducedSS{P}{A_1}) - 1 \le \VCdim(\inducedSS{P}{A_0}) - 1 \le d-1,
\]
where the second inequality holds as $A_1\subset A_0$.
Thus, the new sparsification has dimension at most $d-1$.
Second, the size of the new sparsification is
\[
  s'\coloneqq |B_1|,
\]
and by the choice of $B_1$ we have
\[
  s' \ge c_{\cref{eq:after_sampling}} \cdot |\inducedSS{B_0}{A_0}| \ge c_{\cref{eq:after_sampling}} \cdot  \frac{\alpha}{2 c_{\ref{lem:partition_size}}} |B|.
\]

Finally, let $G'$ be an expansion of $G$ witnessing that the original
$k$-sparsification has complexity $q$,
and let $\phi_1,\ldots,\phi_k$ be formulas defining
$f_1,\ldots,f_k$ in $G'$.
Expand $G'$ further by five unary predicates marking the sets
$A_1$, $A_1'$, $B_1$, $B^*$, and a set $S$ which is all of $V(G)$ if
$\sigma=+$ and empty if $\sigma=-$.
Using the formulas $\phi_1,\ldots,\phi_k$, the relation of belonging to the
same part of $\cal P$ is first-order definable with quantifier rank at most
$q+1$.
Consequently, the edge relation of $M^\sigma$ on $B^*\times A_1$ is also
first-order definable with quantifier rank at most $q+2$:
for $b\in B^*$ and $a\in A_1$, we ask whether there exists
$b'\in B^*$ in the same part of $\cal P$ such that
$N_G(b)\cap (A_1-\set{a})=N_G(b')\cap (A_1-\set{a})$ and
$b,b'$ differ on $a$ in the direction prescribed by $\sigma$.
Restricting this relation to $B_1\times A_1'$ gives a formula defining the
function $f_{k+1}$, of quantifier rank at most $q+2$.
Thus the new sparsification has complexity at most $q+5$.
(We only used quantifier rank $q + 2$, but $q+5$ many unary predicates.)
\end{proof}

\begin{lemma}\label{lem:iterate}
Let $\CC$ be a class in DABE, and let $d$ be its VC-dimension.
Then there is a constant $c_{\ref{lem:iterate}} := c_{\ref{lem:iterate}}(\CC)$ such that every bipartite graph $G=(A,B,E) \in \CC$ of maximal complexity has, for some $i \le d$, a terminal $i$-sparsification  of complexity at most $5d$ and size $s \ge c_{\ref{lem:iterate}} \cdot |B|$.
\end{lemma}
\begin{proof}
  Set $n\coloneqq|A|$ and $m\coloneqq|B|$.
Let $\mathbf S_0$ be the trivial $0$-sparsification of $G$, given by
$A_0\coloneqq A$ and $B_0\coloneqq B$.
Since there are no twins in $B$, this is indeed a sparsification, with
size $s_0=m$, dimension at most $d$, and complexity $c_0=0$.
For $j=0,1,2,\ldots$, construct $\mathbf S_j$ recursively as follows.
Write $s_j$, $d_j$, and $c_j$ for the size, dimension, and
complexity of $\mathbf S_j$, and write $\cal P_j$ for its associated
partition. Finally, let $\alpha_j := s_j/m$

If $\mathbf S_j$ is terminal, or if $j=d$, stop and set $i\coloneqq j$.
Otherwise, \Cref{lem:step} yields a $(j+1)$-sparsification
$\mathbf S_{j+1}$ of size $s_{j+1}$ and complexity $c_{j+1}$ such that
\[
  s_{j+1}
  \ge c_{\ref{lem:step}}(\alpha_j) \cdot m \text{, that is } \alpha_{j+1} \ge c_{\ref{lem:step}}(\alpha_j),
\]
while $c_{j+1}\le c_j+5$ and $d_{j+1}\le d_j-1$.

Let $0 \leq i \leq d$ be the index for which the above construction stops.
Repeated application of the above bounds
gives
\[
  s_i\ge c_{\ref{lem:step}}(c_{\ref{lem:step}}(\dots (1))\dots) \cdot m
  \qquad\text{and}\qquad
  c_i\le 5i\qquad\text{and}\qquad d_i\le d-i.
\]
We argue that $\mathbf S_i$ is terminal,
which will yield the conclusion for $s\coloneqq s_i$ and $c_{\ref{lem:iterate}} := c_{\ref{lem:step}}^d(1)$, where $c_{\ref{lem:step}}$ is seen as a simple function its parameter $\alpha$.

If $i<d$, then $\mathbf S_i$ is terminal by construction.
Suppose now that $i=d$.
Then $\mathbf S_d$ has dimension~$0$.
Let $P\in\cal P_d$.
By definition of the associated partition, the functions
$f^d_1,\ldots,f^d_d$ are constant on $P$.
Since $\mathbf S_d$ is a sparsification, the map
$b\mapsto N_G(b)\cap A_d$ is injective on~$P$. 
Hence, the vertices of $P$ have pairwise different neighborhoods on $A_d$, and we have $|\inducedSS{P}{A_d}| = |P|$.
Furthermore, as $d_i=0$, the set system $\inducedSS{P}{A_d}$ has VC-dimension $0$,
and hence size $1$.
Thus, also $|P|=1$ for every $P\in\cal P_d$, and therefore
\[
  s_d=|B_d|=|\cal P_d|\le 4c_{\ref{lem:partition_size}} \cdot |\cal P_d|.
\]
So $\mathbf S_d$ is terminal as well.
\end{proof}

We now prove our main theorem.

\begin{theorem}\label{thm:main}
Every class in DABE has linear neighborhood complexity.
\end{theorem}
\begin{proof}
Given a class $\CC$, and a graph $G \in \CC$, we let $\text{bip}(G)$ be the bipartite graph $(A, B, E)$ where $A$ and $B$ are two disjoint copies of $V(G)$ and $ab \in E(G)$ when there is an edge between the associated vertices of $V(G)$.
We let $\CC'$ be the hereditary closure of $\{ \text{bip} (G) ~|~ G \in \CC \}$.
Clearly, $\CC'$ is transducible from $\CC$, and $\CC$ has linear neighborhood complexity whenever $\CC'$ has.
Fix now a graph $G=(A,B,E)\in\CC'$, that has maximal relative complexity in $\CC'$ with $n\coloneqq |A| \ge 2$ and $m\coloneqq |B| \ge 1$.
As $\CC'$ is in DABE, the set system $\inducedSS{B}{A}$ has
VC-dimension bounded by a constant $d$ depending only on~$\CC'$ (and so on $\CC$).
By \Cref{lem:iterate}, for some $i\le d$, the graph $G$ has a
terminal $i$-sparsification of size $s$ and complexity at most $5i$ such
that
\[
  s\ge c_{\ref{lem:iterate}} \cdot m.
\]
By \Cref{lem:terminal}, as $i\le d$, we have
\[
  s
  \le c_{\ref{lem:terminal}} \cdot n.
\]
Combining the two inequalities, we obtain
\[
  m
  \le \frac{c_{\ref{lem:terminal}}}{c_{\ref{lem:iterate}}} \cdot n.
\]
Hence, any graph of maximal complexity in $\CC'$ has linearity at most $\frac{c_{\ref{lem:terminal}}}{c_{\ref{lem:iterate}}}$.
By definition, if $\CC'$ contains a graph of linearity $\ell$, then it contains a graph of maximal complexity in $\CC'$ that has linearity at least $\ell$.

\end{proof}

\section{Tiny classes have linear neighborhood complexity}
\begin{lemma}\label{lem:tiny-SP}
Tiny classes are preserved by size-preserving transductions.
\end{lemma}
\begin{proof}
Let $\CC$ be a $c$-tiny class for some constant $c$, and let $T$ be a size-preserving transduction of complexity $k$, with associated interpretation $I$.
In particular, there is a map $T^{-1}: T(\CC) \to \CC$ such that, for each $G \in T(\CC)$, $G \in T(T^{-1}(G))$, and $V(T^{-1}(G)) \le d \cdot |V(G)|$.

Let $n \in \N$ and consider $A_n$ the set of $n$-vertex graphs in $T(\CC)$.
The set $T^{-1}(A_n)$ contains graphs of $\CC$ whose vertex set has size at most $d \cdot n$. In particular, $|T^{-1}(A_n)| \le c^{d \cdot n}$.

Moreover, for each graph $G \in \CC$, a graph of $T(G)$ is obtained by choosing a $k$ unary relations over the $k$-copy of $G$, applying $I$ on this structure (which is a deterministic process), and finally choosing an induced subgraph.
In particular, if $N := |V(G)|$, the number of graphs with $p$ vertices in $T(G)$ is upper bounded by $(2^{kN})^k \cdot \binom{kN}{p}$, where the first term stands for the choice of the $k$ unary relations, and the second for the choice of the induced subgraph.

Therefore, we obtain
\[
|A_n| \le |T^{-1}(A_n)| \cdot \left( 2^{k d\cdot n} \right)^k \cdot \binom{k d \cdot n}{n} \le c^{d \cdot n}2^{k^2 d \cdot n} (ekd)^n \le (\Oh_{k, d, c}(1))^n.
\]

\end{proof}
\begin{theorem}\label{thm:tiny}
Hereditary tiny classes have linear neighborhood complexity
\end{theorem}
\begin{proof}
By [TODO ref], any hereditary non degenerate weakly sparse class of graphs is not tiny. Hence, by \Cref{lem:tiny-SP}, any hereditary tiny class is in SP-DABE.
Applying \Cref{thm:main}, we obtain that hereditary tiny classes have linear neighborhood complexity.
\end{proof}

\section{Substructures of superlinear neighborhood complexity}
\label{sec:obstructions}

We now run the proof of \Cref{thm:main} on a hereditary class $\CC$ of bipartite graphs that is not assumed to be in DABE, but only to have VC-dimension at most $d$ (otherwise $\nc_\CC(n) = 2^n$ for every $n$).
Apart from bounding the VC-dimension, the proof uses DABE only through \Cref{cor:2equiv_ref}, in two places:
\begin{itemize}
  \item in \Cref{lem:terminal}, for the partitions $\Qq_1, \dots, \Qq_k$ of $B_0$ into the preimages of $f_1, \dots, f_k$;
  \item in \Cref{lem:partition_size}, for the associated partition of a sparsification and the partition of $B_0$ into twin classes over some $X \subseteq A_0$ (with $X = A_0$ in \Cref{lem:terminal,lem:step}, and $X$ as in the distillation of \Cref{lem:distil}).
\end{itemize}
Without DABE, each of these applications can fail, giving two partitions whose common refinement is much larger than both of them.
We show that the proof fails only in this way, and describe the partitions involved (\Cref{thm:obstructions}).
For this, we first make explicit how the proof builds its sparsifications.

\subsection{Sparsification steps}
Let $G = (A, B, E)$ be a bipartite graph.
A \remph{sparsification step} turns a $k$-sparsification $(A_k, B_k, f_1, \dots, f_k)$ of $G$, with associated partition $\Pp_{k}$, into a $(k+1)$-sparsification $(A_{k+1} \subseteq A_k,B_{k+1} \subseteq B_k, f_1, \dots, f_{k+1})$ as follows.
\begin{itemize}
  \item Choose $A_{k+1} \subseteq A_k$ and a sign $\sigma_{k+1} \in \set{+, -}$, and let $M_{k+1}$ be the union over $P \in \Pp_k$ of the merge graphs $M^{\sigma_{k+1}}(G[P, A_{k+1}])$ glued along $A_{k+1}$.
  \item Choose $X_{k+1} \subseteq A_{k+1}$ and $B_{k+1} \subseteq B_k$ such that each $b \in B_{k+1}$ has exactly one neighbor in $X_{k+1}$ in $M_{k+1}$, and $(A_{k+1}, B_{k+1}, f_1, \dots, f_k)$ is still a sparsification. Let $f_{k+1}(b)$ be the single neighbor of $b$ in $X_{k+1}$.
\end{itemize}
The associated partition $\Pp_{k+1}$ of the new sparsification is obtained by restricting $\Pp_k$ to $B_{k+1}$ and splitting each part according to $f_{k+1}$.
We only consider sparsifications obtained from a $0$-sparsification $(A_0, B_0)$ by successive sparsification steps, so that $\Pp_0 = \set{B_0}$ and $\Pp_k$ is the partition of $B_k$ into the preimages of $(f_1, \dots, f_k)$.
As in the proof of \Cref{lem:step}, the $k$-sparsification has VC-dimension at most $\VCdim(G) - k$, and complexity $\Oh(k)$.

One can check that \Cref{lem:iterate} obtains all its sparsifications from the trivial $0$-sparsification $(A, B)$ by sparsification steps.

\subsection{Obstructions}

Two partitions $\Pp$ and $\Qq$ of the same set are \remph{$c$-superlinear} if $|\Pp \wedge \Qq| > c \cdot (|\Pp| + |\Qq|)$, that is, if both are small compared to their common refinement.
This is what \Cref{cor:2equiv_ref} forbids for definable partitions in DABE classes.
For a set $Y$, we write $\Pp \res Y$ for the restriction of a partition $\Pp$ to $Y$, and $\twin_X(Y)$ for the partition of $Y$ into twin classes over a set $X$.

\begin{theorem}\label{thm:obstructions}
For all integers $c, d \ge 1$, there is a constant $C = C(c, d)$ such that the following holds.
Let $\CC$ be a hereditary class of bipartite graphs of VC-dimension at most $d$, and let $G = (A, B, E)$ be a graph of maximal complexity in $\CC$.
Then $|B| \le C \cdot |A|$, or $G$ has a sparsification obtained by $k \le d$ sparsification steps with one of the following.
\begin{description}
  \item[(O1) A $c$-superlinear sparsification step:] $|\Pp_k| > c \cdot (|\Pp_{k-1}| + |A|)$.
  \item[(O2) A twin partition and $\Pp_k$ that are $c$-superlinear on some set:] there are $X \subseteq A_k$ and $Y \subseteq B_k$ such that $\Pp_k \res Y$ and $\twin_X(Y)$ are $c$-superlinear.
\end{description}
\end{theorem}

(O1) is where \Cref{lem:terminal} fails: $\Pp_k$ is the common refinement of $\Pp_{k-1} \res {B_k}$ and of the partition of $B_k$ into the preimages of $f_k$, which has at most $|A|$ parts, so these two partitions are $c$-superlinear.
(O2) is where \Cref{lem:partition_size} fails.

\begin{corollary}\label{cor:obstructions}
Let $\CC$ be a hereditary class of graphs of VC-dimension at most $d$ and superlinear neighborhood complexity.
Then for every $c$, some induced subgraph of a graph $\text{bip}(G)$ with $G \in \CC$ has a sparsification obtained by at most $d$ sparsification steps satisfying (O1) or (O2).
\end{corollary}
\begin{proof}
Let $\CC'$ be the hereditary closure of $\setof{\text{bip}(G)}{G \in \CC}$, as in the proof of \Cref{thm:main}.
It has VC-dimension at most $d$ and superlinear neighborhood complexity, so for every $c$ it contains a graph of maximal complexity with $|B| > C(c, d) \cdot |A|$ (see the end of the proof of \Cref{thm:main}), and \Cref{thm:obstructions} applies.
\end{proof}

Since $\Pp_0$ is trivial and the parts of $\Pp_d$ are singletons (the $d$-sparsification has VC-dimension $0$), (O2) requires $1 \le k < d$; as $|\Pp_1| \le |A|$, (O1) requires $k \ge 2$.
In particular, there is no obstruction when $d = 1$, and \Cref{cor:obstructions} recovers that classes of VC-dimension~$1$ have linear neighborhood complexity.

\subsection{Proof of \Cref{thm:obstructions}}

\begin{proof}[Proof of \Cref{thm:obstructions}]
Assume that no sparsification of $G$ obtained by at most $d$ sparsification steps satisfies (O1) or (O2).
We replace the two uses of \Cref{cor:2equiv_ref} in the proof of \Cref{thm:main}.
Consider a sparsification obtained by $k \le d$ sparsification steps.
As (O2) never holds, for every $X \subseteq A_k$,
\begin{equation}\label{eq:no-O2}
  \sum_{P \in \Pp_k} |\inducedSS{P}{X}| \le c \cdot \left(|\Pp_k| + |\inducedSS{B_k}{X}|\right).
\end{equation}
This is the conclusion of \Cref{lem:partition_size}, with $c$ in place of $c_{\ref{lem:partition_size}}$.
As (O1) never holds,
\begin{equation}\label{eq:no-O1}
  |\Pp_k| \le (2c)^d \cdot |A|.
\end{equation}
Indeed, $|\Pp_0| = 1 \le |A|$, and $|\Pp_j| \le c(|\Pp_{j-1}| + |A|) \le (2c)^j |A|$ by induction on $j \le k$.

As the sparsifications of \Cref{lem:iterate} are obtained by sparsification steps, \Cref{lem:distil,lem:step,lem:iterate} apply \Cref{lem:partition_size} only to the associated partitions of such sparsifications, and otherwise only use the bound $d$ on the VC-dimension.
So they hold with \eqref{eq:no-O2} in place of \Cref{lem:partition_size}, with constants depending only on $c$ and $d$: terminality is defined with the threshold $4c$, and the complexity of sparsifications plays no role.
\Cref{lem:iterate} thus gives a terminal sparsification obtained by $k \le d$ sparsification steps, with $|B_k| \ge c_{\ref{lem:iterate}} \cdot |B|$.
As in \Cref{lem:terminal}, using \eqref{eq:no-O2} with $X = A_k$, terminality, and \eqref{eq:no-O1}, we get
\[
  |B_k| = \sum_{P \in \Pp_k} |\inducedSS{P}{A_k}| \le c \cdot \left(|\Pp_k| + |\inducedSS{B_k}{A_k}|\right) \le c (1 + 4c) \cdot |\Pp_k| \le c (1 + 4c) (2c)^d \cdot |A|.
\]
Hence $|B| \le C \cdot |A|$ for $C := c(1 + 4c) (2c)^d / c_{\ref{lem:iterate}}$.
\end{proof}


\section{A transduction of fixed complexity}
\label{sec:fixed}

The obstructions of \Cref{thm:obstructions} are pairs of definable partitions with a superlinear common refinement.
By \Cref{lem:transInter}, a class with obstructions for every $c$ thus transduces their intersection graphs, which form a $2$-degenerate class of superlinear neighborhood complexity.
However, a sparsification obtained by $k$ sparsification steps has complexity $\Oh(k)$, so the complexity of this transduction grows with the VC-dimension.
We now replace the partitions into the preimages of the functions $f_j$ by partitions defined by fixed formulas.

Let $G = (A, B, E)$ be a bipartite graph, $X \subseteq A$, $Y \subseteq B$, $X' \subseteq X$ and $\sigma \in \set{+,-}$.
For $b \in B$, let $D(b)$ be the set of $a \in X'$ such that $N(u) \cap X = (N(b) \cap X) \triangle \set{a}$ for some $u \in Y$, and $a \in N(b)$ if $\sigma = +$, resp.\ $a \notin N(b)$ if $\sigma = -$.
If the partner $u$ were searched in the part of $b$ in $\Pp_{j-1}$ instead of in a fixed set $Y$, then for $(X, X', \sigma) = (A_j, X_j, \sigma_j)$, $D(b)$ would be the set of neighbors of $b$ in $X_j$ in $M_j$, that is, $\set{f_j(b)}$ for $b \in B_j$.
The \remph{simple partition} of $(X, Y, X', \sigma)$ is the partition of the vertices $b$ with $|D(b)| = 1$ according to the element of $D(b)$, and we also call the twin partitions $\twin_X(B)$, for $X \subseteq A$, simple partitions.
Simple partitions are defined by three fixed formulas (one for twin partitions, one for each sign), with at most three unary predicates.
Moreover, these definitions only involve the adjacency between the predicates, so a simple partition of an induced subgraph is the restriction of a simple partition of the whole graph.

\begin{theorem}\label{thm:fixed}
There is a transduction $\mathsf T$ such that, for every hereditary class $\CC$ of graphs of superlinear neighborhood complexity, $\mathsf T(\CC)$ contains a $2$-degenerate class of superlinear neighborhood complexity.
\end{theorem}

The theorem follows from the next lemma, proved in \Cref{sec:proof-simple}.

\begin{lemma}\label{lem:simple}
For all integers $c, d \ge 1$, there is $\kappa = \kappa(c, d)$ such that the following holds.
If a bipartite graph $G = (A, B, E)$ of VC-dimension at most $d$ has a sparsification obtained by at most $d$ sparsification steps satisfying (O1) or (O2) with constant $\kappa$, then there are a set $Y \subseteq B$ and two simple partitions whose restrictions to $Y$ are $c$-superlinear.
\end{lemma}

\begin{proof}[Proof of \Cref{thm:fixed}]
We first show that for every $c$, some graph $\text{bip}(G)$ with $G \in \CC$ has a set $Y \subseteq B$ and two simple partitions whose restrictions to $Y$ are $c$-superlinear.
If $\CC$ has VC-dimension at most $d$, this follows from \Cref{cor:obstructions}, applied with $\kappa(c, d)$ in place of $c$, \Cref{lem:simple}, and the above remark on induced subgraphs.
Otherwise, for every $n$ some $G \in \CC$ has a shattered set of size $2n$, whose copy in $A$ in $\text{bip}(G)$ we split into $X_1 \uplus X_2$ with $|X_1| = |X_2| = n$.
Let $Y \subseteq B$ contain, for all $x_1 \in X_1$ and $x_2 \in X_2$, a vertex $b$ with $N(b) \cap (X_1 \cup X_2) = \set{x_1, x_2}$.
Then $\twin_{X_1}(Y)$ and $\twin_{X_2}(Y)$ have $n$ parts each, and their common refinement has $n^2$ parts, so they are $c$-superlinear when $n > 2c$.

As simple partitions are defined by finitely many formulas, \Cref{lem:transInter,lem:transunion} give a transduction $\mathsf T'$ such that, for every bipartite graph $H$, $\mathsf T'(H)$ contains the intersection graphs $\text{Inter}(\Pp, \Qq)$ of the restrictions $\Pp, \Qq$ of any two simple partitions of $H$ to any set $Y \subseteq B$.
These graphs are $2$-degenerate: the vertices of $\Pp \wedge \Qq$ have degree $2$, and the other vertices are pairwise non-adjacent.
For the $c$-superlinear pairs found above, the more than $c \cdot (|\Pp| + |\Qq|)$ vertices of $\Pp \wedge \Qq$ have pairwise distinct neighborhoods in $\Pp \uplus \Qq$.
As $c$ is arbitrary, the theorem holds with $\mathsf T := \mathsf T' \circ \text{bip}$.
\end{proof}

\subsection{Proof of \Cref{lem:simple}}
\label{sec:proof-simple}

Two partitions of the same set \remph{intersect $c$-linearly} if their restrictions to any subset are not $c$-superlinear.

\begin{lemma}\label{lem:3-partitions}
Let $\Pp, \Qq, \mathcal R$ be partitions of a set $S$ such that $\Pp \wedge \mathcal R$ refines $\Pp \wedge \Qq$, and $\Pp$ and $\mathcal R$ intersect $c$-linearly for some $c \ge 1$.
If $|\Pp \wedge \Qq| \ge 4c|\Pp|$, then some $S' \subseteq S$ containing at most one element of each part of $\mathcal R$ meets at least $|\Pp \wedge \Qq| / (4c)$ parts of $\Pp \wedge \Qq$.
\end{lemma}
\begin{proof}
Let $\mathcal L$ be the set of parts of $\mathcal R$ meeting at least $2c$ parts of $\Pp$, let $L := \bigcup \mathcal L$, and let $\Pp \wedge \mathcal L$ be the set of nonempty sets $P \cap R$ with $P \in \Pp$ and $R \in \mathcal L$.
As the restrictions of $\Pp$ and $\mathcal R$ to $L$ are not $c$-superlinear, $2c|\mathcal L| \le |\Pp \wedge \mathcal L| \le c(|\Pp| + |\mathcal L|)$.
Hence $|\mathcal L| \le |\Pp|$ and $|\Pp \wedge \mathcal L| \le 2c|\Pp|$.
Each part of $\Pp \wedge \Qq$ is a union of parts of $\Pp \wedge \mathcal R$, and if it meets $L$, it contains a part of $\Pp \wedge \mathcal L$.
So at most $2c|\Pp| \le |\Pp \wedge \Qq|/2$ parts of $\Pp \wedge \Qq$ meet $L$.

For every part $R \notin \mathcal L$ of $\mathcal R$, choose $x_R$ uniformly at random in a set of fewer than $2c$ elements of $R$ meeting every nonempty set $P \cap R$ with $P \in \Pp$.
A part of $\Pp \wedge \Qq$ disjoint from $L$ contains a nonempty set $P \cap R$ with $R \notin \mathcal L$, so it contains $x_R$ with probability more than $1/(2c)$.
Hence $S' := \setof{x_R}{R \in \mathcal R \setminus \mathcal L}$ meets in expectation more than $|\Pp \wedge \Qq|/(4c)$ parts of $\Pp \wedge \Qq$, and we fix such an outcome.
\end{proof}

We use \Cref{lem:3-partitions} with $\mathcal R = \twin_{A_t}(S)$, to make the ground set twin-free over $A_t$.
On such a set, the partitions into the preimages of the $f_j$ are simple partitions, up to a constant loss.

\begin{lemma}\label{lem:flatten}
Let $G = (A, B, E)$ be a bipartite graph of VC-dimension at most $d$ with a sparsification obtained by $t$ sparsification steps, and let $Y \subseteq B_t$ be twin-free over $A_t$.
There are $Y' \subseteq Y$ with $|Y'| \ge 2^{-4dt}|Y|$ and simple partitions $\mathcal N_1, \dots, \mathcal N_t$ such that, for every $j \le t$, the restriction of $\mathcal N_j$ to $Y'$ is the partition of $Y'$ into the preimages of $f_j$.
\end{lemma}
\begin{proof}
Let $j \le t$.
For $b \in Y$, let $p_j(b)$ be the partner of $b$ that gives its edge to $f_j(b)$ in $M_j$, that is, a vertex in the part of $b$ in $\Pp_{j-1}$ with $N(p_j(b)) \cap A_j = (N(b) \cap A_j) \triangle \set{f_j(b)}$.
Let $Y_j := p_j(Y)$, and let $D_j(b)$ be the set $D(b)$ defined for $(A_j, Y_j, X_j, \sigma_j)$; it contains $f_j(b)$.
As $A_t \subseteq A_j$, $Y$ is also twin-free over $A_j$.
So $\sum_{b \in Y} |D_j(b)|$ is at most the number of ordered pairs of sets at Hamming distance $1$ in the set system of the traces on $A_j$ of the vertices of $Y \cup Y_j$, which is at most $2d(|Y| + |Y_j|) \le 4d|Y|$ by \Cref{lem:VC-dim2deg}.

For each $j \le t$ independently, let $X'_j$ be a random subset of $X_j$ containing each element with probability $1/2$.
A vertex $b \in Y$ satisfies $D_j(b) \cap X'_j = \set{f_j(b)}$ for all $j \le t$ with probability $2^{-\sum_j |D_j(b)|}$, so by convexity the expected number of such vertices is at least $2^{-4dt}|Y|$.
Fix such an outcome, let $Y'$ be the set of these vertices, and let $\mathcal N_j$ be the simple partition of $(A_j, Y_j, X'_j, \sigma_j)$.
For $b \in Y'$, the set $D(b)$ defined for $(A_j, Y_j, X'_j, \sigma_j)$ is $D_j(b) \cap X'_j = \set{f_j(b)}$, so $b$ lies in the part of $\mathcal N_j$ associated with $f_j(b)$.
\end{proof}

Finally, $\Pp_t$ is the partition into the preimages of $(f_1, \dots, f_t)$, and the next lemma extracts a superlinear pair from such a partition.
For sets $X_1, \dots, X_k$, a set $R \subseteq X_1 \times \dots \times X_k$ and $i < j$, let $\pi_{i,j}(R) := \setof{(x_i, x_j)}{(x_1, \dots, x_k) \in R}$.

\begin{lemma}\label{lem:projection}
Let $c, k \ge 1$, let $X_1, \dots, X_k$ be finite sets, and let $R \subseteq X_1 \times \dots \times X_k$ with $|R| > 2^{2ck} \cdot (|X_1| + \dots + |X_k|)$.
Then there are $i < j$, $X'_i \subseteq X_i$ and $X'_j \subseteq X_j$ such that $|\pi_{i,j}(R) \cap (X'_i \times X'_j)| > c \cdot (|X'_i| + |X'_j|)$.
\end{lemma}
\begin{proof}
We may assume that the sets $X_i$ are pairwise disjoint.
Let $H$ be the graph on $X_1 \cup \dots \cup X_k$ whose edges are the pairs $x_i x_j$ with $(x_i, x_j) \in \pi_{i,j}(R)$ for some $i < j$.
If the lemma fails, then every $U \subseteq V(H)$ satisfies, with $X'_i := U \cap X_i$,
\[
  |E(H[U])| = \sum_{i < j} |\pi_{i,j}(R) \cap (X'_i \times X'_j)| \le c \sum_{i < j} (|X'_i| + |X'_j|) = c (k - 1) \cdot |U|.
\]
Hence $H$ is $D$-degenerate for $D := 2c(k-1)$.
In a degeneracy order, a clique of $H$ is determined by its first vertex and a subset of the at most $D$ later neighbors of this vertex, so $H$ has at most $2^D \cdot |V(H)|$ cliques.
Every tuple of $R$ spans a clique of $H$, and distinct tuples span distinct cliques, so $|R| \le 2^D \cdot |V(H)|$, a contradiction.
\end{proof}

\begin{proof}[Proof of \Cref{lem:simple}]
Let $\kappa_0 := 1$, let $\kappa_t := 4t\kappa_{t-1} \cdot 2^{4dt + 2c(t+1)}$ for $t \ge 1$, and let $\kappa := \kappa_d$.
Let $t$ be minimal such that $G$ has a sparsification obtained by $t$ sparsification steps satisfying (O1) or (O2) with constant $\kappa_t$; it exists as the $\kappa_t$ increase, and $t \ge 1$ as $\Pp_0$ is trivial.
By minimality, the $(t-1)$-sparsification from which it is obtained does not satisfy (O2) with constant $\kappa_{t-1}$, with $A_t$ in place of $X$ and any subset of $B_t$ in place of $Y$.
That is, $\Pp_{t-1} \res {B_t}$ and $\twin_{A_t}(B_t)$ intersect $\kappa_{t-1}$-linearly.

In case (O1), let $Y := B_t$ and $\Qq := \Pp_t$.
In case (O2), let $X$ and $Y$ be as in (O2), and $\Qq := (\Pp_t \res Y) \wedge \twin_X(Y)$.
In both cases, $\Qq$ refines $\Pp_{t-1} \res Y$, and $|\Qq| > \kappa_t |\Pp_{t-1} \res Y| \ge 4\kappa_{t-1} |\Pp_{t-1} \res Y|$.
As $(A_t, B_t, f_1, \dots, f_{t-1})$ is a sparsification, the common refinement of $\Pp_{t-1} \res Y$ and $\twin_{A_t}(Y)$ is the partition of $Y$ into singletons.
So \Cref{lem:3-partitions}, applied on $Y$ to $\Pp_{t-1} \res Y$, $\Qq$ and $\twin_{A_t}(Y)$, gives a set $Y_1 \subseteq Y$ that is twin-free over $A_t$ and meets at least $|\Qq|/(4\kappa_{t-1})$ parts of $\Qq$; we keep in $Y_1$ one vertex of each of these parts.

The partition $\Qq$ of $Y$ is into the preimages of the map $b \mapsto (f_1(b), \dots, f_t(b))$, extended by $N(b) \cap X$ in case (O2), so this map is injective on $Y_1$.
Let $\Sigma$ be the sum of the sizes of the images of these $s \le t + 1$ coordinates.
In case (O1), $\Sigma \le t|A| < t|\Qq|/\kappa_t$.
In case (O2), as $\Pp_t$ refines the partition into the preimages of each $f_j$, $\Sigma \le t|\Pp_t \res Y| + |\inducedSS{Y}{X}| < t|\Qq|/\kappa_t$.
\Cref{lem:flatten} gives $Y_2 \subseteq Y_1$ and simple partitions $\mathcal N_1, \dots, \mathcal N_t$ whose restrictions to $Y_2$ are the partitions into the preimages of $f_1, \dots, f_t$, with
\[
  |Y_2| \ge 2^{-4dt}|Y_1| \ge \frac{2^{-4dt}}{4\kappa_{t-1}} |\Qq| > \frac{2^{-4dt} \kappa_t}{4t\kappa_{t-1}} \Sigma = 2^{2c(t+1)} \Sigma.
\]
By \Cref{lem:projection}, applied to the image of $Y_2$, there are two coordinates $i < j$ and subsets $X'_i, X'_j$ of their images such that on the set $Y_3$ of vertices of $Y_2$ whose $i$-th and $j$-th coordinates lie in $X'_i$ and $X'_j$, the partitions according to these two coordinates have at most $|X'_i|$ and $|X'_j|$ parts, while their common refinement has more than $c(|X'_i| + |X'_j|)$ parts.
These two partitions are the restrictions to $Y_3$ of simple partitions (of the form $\mathcal N_\ell$ or $\twin_X(B)$).
\end{proof}

\end{document}